\paragraph{Cash classification is a supplied premise, not a balance-sheet inference.}
\label{sec:credit-cash-classification}
The separate \texttt{CreditCashClassification/1} layer makes that missing
premise explicit. It imports an exact retained credit projection and asks the
caller to identify the classification policy and its context. Each price,
integral cost, integral fee and collection retains its native date, amount and
journal origin. Effective-interest accrual and impairment changes are not
cash flows merely because they affect income.

Consider the retained example with purchase price 98, integral cost 2, zero
integral fee, and subsequent collections of 44 and 55. For a purely synthetic
classification test, suppose the caller assigns price and collections to
investing and cost and fee to operating. Investing then has receipts of 99,
payments of 98 and net cash of 1; operating has payments of 2 and net cash of
$-2$. Together they explain the cash reduction of 1. The example checks the
consequences of those supplied labels. It does not select an IAS 7 policy for
an actual lender or establish which part of a receipt is contractual interest.

A classification names a native origin by its position and content digest.
The caller can declare separate components, but their exact signed amounts
must sum to that origin. Positive and negative components cannot cancel within
one supposedly gross receipt or payment. A principal/interest split therefore
needs explicit supplied amounts and supporting policy; it cannot be copied
from the effective-interest accrual. The original gross origin and each declared
component remain available alongside the policy reference.

The classification follows journal prefixes, not dates alone. At each native
observation it includes only cash origins already introduced by that journal
prefix and reports gross receipts, gross payments and net amounts by activity.
It preserves the exact financial balances and profits from the original credit
projection. A changed activity label can change the cash-flow presentation but
cannot change the native journal or the closing cash balance.

Missing classification remains unresolved even when its amount is zero. In the
example, omitting the zero integral fee leaves an incomplete origin population;
the unclassified net is zero, but that numerical equality does not prove complete
classification. Likewise, an unknown later collection leaves earlier completed
prefixes usable while making the affected later statement partial. Known
activity totals remain labelled as the classified subset until every origin
has a declaration. There is no balancing residual category.

For a rational financial state, the adapter supplies these explicit flows to
the unchanged \texttt{CompanyBookReport/1}. Across the retained nine schedules,
two account charts and presence or absence of the unrelated equipment balance,
the three conditional classification examples produce 684 observations.
The 372 rational observations agree with the host cash, balance, income and
equity reports. The remaining 312 keep their single-root financial expressions;
the adapter does not round them to obtain host-book acceptance. Counterexamples
cover changed source and policy bindings, duplicate origins, wrong signs or
totals, unsupported activities and concealed netting.

This supplies conditional classified reporting for the retained projections,
while preserving the original unclassified report's refusal. Actual issuer
facts and classification-policy eligibility, complete notes, fiscal completeness,
the full indirect cash bridge and arbitrary multi-root contracts remain separate
requirements. Exact cases, unknowns and refusals are retained under
{\small\path{data/accounting-master/credit-cash-classification-20261026/}}.

